% Packaging the codebase as a wheel, and where its time goes (build/packaging).
\documentclass[11pt]{article}
\usepackage[T1]{fontenc}
\usepackage{lmodern}
\usepackage{microtype}
\usepackage[margin=1in]{geometry}
\usepackage{amsmath,amssymb,amsthm,mathtools}
\usepackage{algorithm}
\usepackage{algpseudocode}
\usepackage{booktabs}
\usepackage[numbers,sort&compress]{natbib}
\usepackage[hidelinks]{hyperref}
\usepackage[capitalise,noabbrev]{cleveref}

\theoremstyle{plain}
\newtheorem{theorem}{Theorem}[section]
\newtheorem{lemma}[theorem]{Lemma}
\newtheorem{proposition}[theorem]{Proposition}
\newtheorem{corollary}[theorem]{Corollary}
\theoremstyle{definition}
\newtheorem{definition}[theorem]{Definition}
\theoremstyle{remark}
\newtheorem{remark}[theorem]{Remark}

\newcommand{\lean}[1]{\texttt{#1}}
\newcommand{\envroot}{\texttt{MIND\_PALACE\_ROOT}}

\title{Packaging a Research Codebase and Measuring Where Its Time Goes}
\author{mind-palace contributors}
\date{7 November 2026}

\begin{document}
\maketitle

\begin{abstract}
We package a research knowledge-store codebase as a pure-Python wheel and measure where its time
goes before rewriting any of it in a compiled language. Packaging needs one rule for the instance
root (an environment variable, else the checkout, else a refusal at first use), a console script,
and a CI job that installs the wheel away from the checkout. A profile that controls for the
profiler's blind spots puts over 97\% of a query in building its universe, mostly LanceDB and
SQLite, and the interpreter at 0.304\,\textmu s a step. Five changes kept in Python take a query
from 93 to about 20\,ms, a walk solve from 7.3\,s to 286\,ms and a reachability test from 36.5\,s to
82\,ms; no operation argues for a compiled rewrite.
\end{abstract}

\section{Introduction}

The mind-palace codebase is a knowledge store with a query algebra: a sentence such as
\texttt{near "\dots"} or \texttt{reach F from <note>} is compiled, planned, evaluated over a
universe of nodes gathered from a vector store (LanceDB) and a citation graph (SQLite), and ranked.
It also holds a small interpreter for scripts, a statistical scorer, and a temporal index over the
stores' events. Until this change it ran only from a git checkout, and six of its operations had
been named as candidates for a compiled extension (a Rust crate through PyO3).

This paper does two things that belong together. It makes the codebase installable, which any
compiled extension would also need, and it measures where the time goes in the six candidate
operations before deciding whether any of them should leave Python. The contributions are:
\begin{itemize}
\item One rule for the instance root, a refusal at first use rather than at import, a
  \texttt{mind-palace} console script, and a CI job that builds the wheel, installs it in a fresh
  environment away from the checkout and runs it (\cref{sec:root,sec:impl}).
\item A profiling method for a pipeline that mixes Python with compiled libraries, with a control
  for each blind spot of a deterministic profiler, and the resulting profile (\cref{sec:profile}).
\item Five changes kept in Python, each measured before and after in one session and each pinned
  by a test against the code it replaced: reading HEAD from the repository's files, reading vectors
  through Arrow, batching the citation graph's reads, solving the killed walk on the source's
  connected component, and a bitset reachability index for the temporal spine
  (\cref{sec:model,sec:changes}).
\item A verdict per candidate, bounded by Amdahl's law from the measured share of time spent in
  Python, with the costs a compiled rewrite would carry (\cref{sec:verdict}), and an independent
  re-measurement of six figures in a second session (\cref{sec:remeasure}).
\end{itemize}

\section{Background and related work}

\paragraph{Measure first.} \citet{knuth1974} warned that programmers spend much of their effort
on the speed of noncritical parts and that the critical few percent should be found by measurement,
not intuition. \citet{gregg2020} sets out methodologies for performance analysis that start from
the workload and its observed costs rather than from a suspected component. \citet{amdahl1967}
bounded the gain from speeding up part of a computation by the share of time that part takes, the
bound we apply to each candidate. \citet{mytkowicz2009} showed that innocuous details of an
experimental setup can bias a performance comparison, one reason we measure each change against its
predecessor back to back in one session and report medians.

\paragraph{The compiled layers.} Most of the codebase's time is already spent in compiled code:
NumPy~\cite{harris2020} and SciPy~\cite{virtanen2020} for the numeric kernels, SQLite for the
citation graph~\cite{gaffney2022}, Apache Arrow~\cite{arrow} as LanceDB's~\cite{lancedb} columnar
format. Python's deterministic profiler, cProfile~\cite{pyprofile}, records Python frames and calls
to built-in functions; time inside a NumPy operator, a LAPACK routine or Arrow's Cython code lands
in the calling frame's own time, so a naive reading overstates Python.

\paragraph{Reachability.} The temporal spine is a directed acyclic graph of events, and its
restriction to one stratum needs the transitive closure among that stratum's events.
\citet{purdom1970} computes the closure through strongly connected components and a topological
order, merging successor sets as bit vectors; \citet{nuutila1995} improves this with interval
representations of the closure; \citet{agrawal1989} compress the closure of a graph by labelling a
spanning tree with intervals and propagating non-tree edges. The spine is a union of chains joined
by cross-store edges whose heads have several producers, a DAG and not a forest, so interval labels
degrade to propagated sets and the bit-vector approach is the direct one.

\paragraph{The walk.} The \texttt{diffuse} sentence ranks nodes by the free-energy distance of
randomized shortest paths~\cite{kivimaki2014}, computed from the inverse of $I - W_\beta$ for a
killed walk $W_\beta$ on the support graph.

\paragraph{Packaging.} A wheel is Python's built distribution format~\cite{pep427}, and
\texttt{pyproject.toml} names the build backend that makes one~\cite{pep517}; a wheel tagged
\texttt{py3-none-any} is pure Python and installs on any platform.

\section{Model}\label{sec:model}

No Lean declaration checks the statements below. Each is checked by tests against the code it
replaced (\cref{sec:tests}).

\subsection{The instance root}\label{sec:root}

\begin{definition}[Root rule]
Let $E$ be the environment and $c$ the checkout the code runs from, or $\bot$ for an installed copy
(a checkout is a code directory holding \texttt{pyproject.toml} and \texttt{core/kernel}). The root
is $E[\envroot]$ when that is set, and must then name an existing directory or the rule refuses;
otherwise it is $c$; otherwise there is no root, and the first filesystem use of a path under it
raises an error naming the variable.
\end{definition}

The root holds what an instance owns and the code does not ship: its data directory, its
per-machine overlays and its constitution. The committed defaults and the script library ship
inside the wheel and are found beside the code, so they need no root. Two properties follow from the
definition and are pinned by tests: in a checkout with the variable unset, the root and every path
derived from it equal what they were before packaging; and importing the packages needs no root,
because the unresolved root is a path object that joins like any path but raises on any file system
use, so \texttt{import core} works anywhere (the CI job checks it) and a stray join cannot land in
the working directory.

\subsection{What a rewrite can buy}

\begin{proposition}[Amdahl bound]\label{prop:amdahl}
If a fraction $f$ of an operation's wall time is spent in code that a rewrite speeds up by a factor
$s$, the operation's speedup is $1/((1 - f) + f/s) < 1/(1 - f)$.
\end{proposition}

We call $f$ the operation's \emph{Python share}: the measured fraction of wall time that is not
compiled code. A compiled extension that made the Python free would give at most $1/(1-f)$: 1.25
for $f = 0.2$ and 2 for $f = 0.5$. An operation clears a fivefold end-to-end gain only when $f >
0.8$.

\subsection{The killed walk decomposes by component}

Let $A$ be the symmetric non-negative weight matrix of the support graph, $D$ a diagonal matrix of
degrees (the support's own, or the full graph's for a support whose complement absorbs), $P = D^{-1}
A$, and $W_\beta = z\, P \circ A^{\circ \beta}$ with $0 < z < 1$ and $\beta > 0$, the powers taken
entrywise. The free-energy distance uses $Z = (I - W_\beta)^{-1}$ through $\varphi_{ij} = -\log(Z_{ij} /
Z_{jj}) / \beta$, symmetrized as $(\varphi_{ij} + \varphi_{ji})/2$, with $+\infty$ where $Z_{ij} =
0$.

\begin{proposition}[Block structure]\label{prop:block}
Order the nodes by connected component $C_1, \dots, C_r$. Then $W_\beta$ and $Z$ are block diagonal
with blocks $C_1, \dots, C_r$, and for a source $s \in C_q$ the distances from $s$ are those
computed from $(I - W_{\beta, q})^{-1}$ on $C_q$ alone, and $+\infty$ to every node outside $C_q$.
\end{proposition}
\begin{proof}
$A_{ij} = 0$ whenever $i$ and $j$ lie in different components. Row scaling by $D^{-1}$, the entrywise
power and the entrywise product preserve this zero pattern, so $W_\beta = \mathrm{diag}(W_{\beta,1},
\dots, W_{\beta,r})$, and the inverse of a block-diagonal matrix is the block-diagonal matrix of the
blocks' inverses. Every quantity the distance from $s$ reads ($Z_{sj}$, $Z_{js}$, $Z_{jj}$ for $j
\in C_q$) lies in block $q$, and $Z_{sj} = 0$ for $j \notin C_q$.
\end{proof}

The full inverse costs $O(n^3)$ time and $O(n^2)$ memory for $n$ nodes; the block costs
$O(|C_q|^3)$ and $O(|C_q|^2)$, plus $O(n + m)$ to find the components. On a connected support the
two coincide and the component search is pure overhead.

\subsection{The reachability index}

\begin{proposition}[Closure by bitsets]\label{prop:bits}
Let $G = (V, E)$ be a DAG with $n$ nodes and $m$ edges, numbered in a topological order $\tau$, and
define $\mathrm{desc}(u) = \bigcup_{(u,v) \in E} (\{v\} \cup \mathrm{desc}(v))$, computed in reverse
order of $\tau$ as an $n$-bit integer. Then $\mathrm{desc}(u)$ is the set of strict descendants of
$u$, its computation takes $O(m\, n / w)$ word operations for words of $w$ bits, and for a set $K$
of kept nodes with indicator bitset $\chi_K$, $\mathrm{desc}(u) \wedge \chi_K$ is the set of strict
descendants of $u$ in $K$.
\end{proposition}
\begin{proof}[Proof sketch]
Every successor of $u$ follows $u$ in $\tau$, so its set is complete when $u$ is processed; the
claim follows by induction in reverse order. Each edge costs one OR of two $n$-bit integers. A
descendant of $u$ in $K$ is a descendant of $u$ that is in $K$.
\end{proof}

A restriction of the spine to a stratum $K$ is then the mask $\chi_K$ over a single index built once
per spine, and the closed adjacency of a kept event is decoded only when it is read. The closure it
replaced ran a breadth-first search per kept event, $O(|K|\,(n + m))$, and stored the result.

\subsection{What the other changes keep}

The HEAD reader returns either the commit \texttt{git rev-parse HEAD} would, or nothing, in which
case the caller asks git as before; it resolves a \texttt{.git} directory or a worktree's
\texttt{gitdir} file, the common directory of a linked worktree, a loose or a packed ref, and a
detached HEAD. The vector read widens half- and single-precision values to double, which is exact,
so the lists it returns are bit for bit those of the row-by-row conversion. The batched citation
read returns the set of endpoint pairs that the per-node reads of both directions return.

\section{Concepts and intuition}

\paragraph{The profiler sees Python.} A deterministic profiler attributes time to Python frames.
Time spent inside compiled code appears as the caller's own time, and time spent waiting for another
runtime appears as a wait in an event loop or a poll. Reading such a profile literally makes every
pipeline look like Python. The method below assigns each of those to the compiled side by a control:
the same call at the smallest size that exercises the same code path measures the interpreter's
share, and waits are attributed by their callers.

\paragraph{A missing cache is not a slow language.} The largest single frame in a stopped
instance's queries was two git child processes asked for a commit that the repository's files
already hold. Removing a needless call is a change of design, not of language.

\paragraph{Ask the compiled layer for the shape you need.} LanceDB returned vectors as Python lists
of floats built row by row; asking for the Arrow column and converting it with NumPy in one step
took the same values at a fraction of the cost. SQLite answered two statements per node with full
rows when the builder needed two columns for a set of nodes.

\paragraph{Structure before speed.} The walk solve was already LAPACK, so no rewrite could speed it
up; its matrix's block structure lets it solve a smaller problem. The spine's restriction was all
Python and quadratic; a different representation of the closure made it two orders of magnitude
faster without leaving Python. Both kept the code they replaced as an oracle in the tests.

\section{Implementation}\label{sec:impl}

\paragraph{The wheel.} \texttt{pyproject.toml} names hatchling as the build backend and takes the
version from the git tag through hatch-vcs, with \texttt{0.0.0+unknown} as the fallback for a build
without a repository. The importable packages are listed rather than discovered (agents, config,
core, edge, eval, ops, scheduler, scripts); the script library ships beside them as data; the
per-machine overlays are excluded. The \texttt{mind-palace} console script is the same verb table
as the checkout's shell front door, and a test pins the two equal. The CI job builds the wheel,
creates a fresh environment, installs the wheel into it, changes to a directory away from the
checkout, imports three packages and runs two commands.

\begin{algorithm}[t]
\caption{Restricting the spine to a stratum through the reachability index.}\label{alg:index}
\begin{algorithmic}[1]
\Function{Build}{$\mathit{succ}$} \Comment{once per spine, cached on the instance}
  \State $\tau \gets$ Kahn's topological order of $\mathit{succ}$; \quad $\mathit{pos}[\tau_i] \gets i$
  \For{$i \gets n-1$ \textbf{down to} $0$}
    \State $\mathit{desc}[i] \gets \bigvee_{v \in \mathit{succ}(\tau_i)}
      \bigl(\mathit{desc}[\mathit{pos}[v]] \vee 2^{\mathit{pos}[v]}\bigr)$
  \EndFor
  \State \Return $(\tau, \mathit{pos}, \mathit{desc})$
\EndFunction
\Function{Restrict}{$\mathit{spine}, s$}
  \State $K \gets$ the events of stratum $s$; \quad $\chi \gets$ the bitset of $K$, built through
    bytes
  \State \Return a spine on $K$ whose closed adjacency of $e \in K$ is decoded from
    $\mathit{desc}[\mathit{pos}[e]] \wedge \chi$ when read
\EndFunction
\end{algorithmic}
\end{algorithm}

\begin{algorithm}[t]
\caption{One row of the free-energy distance.}\label{alg:component}
\begin{algorithmic}[1]
\Require weights $A$ ($n \times n$), source $s$, $\beta$, $z$, optional degrees
\If{$n \le 300$} \Return row $s$ of the full computation \EndIf
\State $\mathit{label} \gets$ the connected components of $A$ (sparse, undirected)
\State $C \gets \{j : \mathit{label}[j] = \mathit{label}[s]\}$
\State $\mathit{out} \gets (+\infty)^n$
\State $\mathit{out}[C] \gets$ row of $s$ in the full computation on $A[C, C]$ with degrees
  restricted to $C$
\State \Return $\mathit{out}$
\end{algorithmic}
\end{algorithm}

\Cref{alg:index} is the spine's index. Python's arbitrary-precision integers serve as bitsets, so an
OR of two $n$-bit sets is a single operation over $n/64$ machine words, and decoding unpacks the
bytes with NumPy. The index is a cached property, which writes the instance dictionary and so holds
on a frozen dataclass; a restriction of a restriction masks the parent's bitsets instead of building
a second closure. \Cref{alg:component} is the walk's row; the route changes above 300 nodes. A
sparse direct solve was built and measured first and
lost: the distance needs the diagonal of the inverse, one solve per node, and fill makes each solve
expensive, so on a connected support of 5{,}000 nodes it took 34\,s against the dense inverse's
1.2\,s.

The other three changes are local. The HEAD reader replaces the call that ran \texttt{git
rev-parse HEAD} and \texttt{git status}, whose dirty flag the view never read; anything it cannot
place goes to git as before. The vector read asks LanceDB for an Arrow table, combines the vector
column, reshapes its values into one NumPy matrix, widens it to double and splits it into rows. The
citation store answers \texttt{neighbour\_pairs(refs, commit)}: the two endpoint columns of the
edges leaving or entering a set of refs, one statement per 400 refs per direction, and the universe
builder's search fetches its frontier ahead with it.

\section{Evaluation}\label{sec:eval}

\subsection{Method}\label{sec:method}
The profile ran on synthetic fixtures at the live sizes, built in a scratch directory and never on
the live stores: 17{,}328 vectors of 768 half-precision dimensions in LanceDB, of which 304 are the
chunks of 38 notes; 147{,}057 curated atoms; 3{,}029{,}562 citation edges over 1{,}487 commits in
SQLite with the real schema and indexes and a skewed degree distribution whose hub has 684 citers;
and a fake embedder, so no model time enters any figure. Wall times are medians of 3 to 7
unprofiled runs on a warm page cache, and the split between components comes from one cProfile run
of the same call, with three corrections: NumPy operators, LAPACK and Arrow's Cython are invisible
to the profiler, so a numeric kernel's Python share is measured by the floor control (the same call
at the smallest size that exercises its code path); LanceDB's synchronous calls wait in the event
loop's \texttt{kqueue} and a git child process waits in \texttt{poll}, and each wait is attributed by
its caller; and LanceDB's list conversion is Arrow's Cython, counted as compiled. Each kept change
was measured against its predecessor back to back in one session, the two arms interleaved, as
medians; memory was measured in separate processes as peak resident size. The machine was a laptop
with the instance's own daemon running, which adds a few percent of noise. The profiling harness is
not committed; its parameters are those stated here, and \cref{sec:remeasure} repeats six of its
figures with scripts described in full.

\subsection{Where the time goes}\label{sec:profile}
In every sentence, compilation takes 0.1 to 0.5\,ms, evaluation 0.1 to 1.0\,ms and ranking at most
0.8\,ms; the universe build is over 97\% of the wall time. \Cref{tab:pipeline} splits it. The
\texttt{near} family is LanceDB's, and its largest frame was not the search but the conversion of
the 304 note vectors into Python lists. The citation-graph family is half SQLite (two statements
per node) and half Python object construction. With no run active, every sentence over the citation
graph also started two git processes, 25 to 45\,ms of it. Elsewhere, the killed-walk solve is
LAPACK (under 1\% Python at 1{,}000 nodes and above), the ingress interpreter is entirely Python at
0.304\,\textmu s a step (0.299 to 0.312 across 5{,}000 to 120{,}000 steps, a quarter of it the
meter's own charging), the scorer's bootstrap is 98\% NumPy, and the spine's partition test was
entirely Python and quadratic, 4.2 to 4.3 times slower per doubling of the events.

\begin{table}[ht]
\centering\small
\caption{A sentence's support, its wall time with a run active, and the split of its time:
LanceDB with Arrow, SQLite, NumPy with SciPy, and Python (profile of the base
commit).}\label{tab:pipeline}
\begin{tabular}{@{}lrrrrrr@{}}
\toprule
Sentence & Nodes & Wall & LanceDB & SQLite & NumPy & Python \\
\midrule
\texttt{near "\dots"} & 38 & 51\,ms & 73\% & 2\% & 7\% & 17\% \\
\texttt{near \dots\ over curated} & 17 & 78\,ms & 84\% & 1\% & 3\% & 10\% \\
\texttt{diffuse S from <note>} & 44 & 58\,ms & 68\% & 6\% & 7\% & 18\% \\
\texttt{walk from <note> via F\{1,3\}} & 200 & 26\,ms & 0\% & 49\% & 0\% & 50\% \\
\texttt{reach F from <note>} & 7 & 5.9\,ms & 0\% & 46\% & 0\% & 54\% \\
\texttt{reach F from <hub file>} & 200 & 14.6\,ms & 0\% & 45\% & 0\% & 54\% \\
\bottomrule
\end{tabular}
\end{table}

\subsection{The five kept changes}\label{sec:changes}
\Cref{tab:changes} gives the three changes on the query path; stacked, they take \texttt{near} from
93 to about 20\,ms and \texttt{reach F} from a note from 42 to about 2\,ms. The cold first query
over the curated store (1.43\,s, 84\% inside LanceDB) is not on the vector-read path and did not
move. \Cref{tab:walk} gives the walk solve, and \cref{tab:spine} the spine's index.

\begin{table}[ht]
\centering\small
\caption{The query path before and after, A/B in one session, interleaved, medians.}
\label{tab:changes}
\begin{tabular}{@{}llrr@{}}
\toprule
Change & Sentence & Before & After \\
\midrule
HEAD read from \texttt{.git} (no run active) & \texttt{near} & 93.1\,ms & 55.4\,ms \\
& \texttt{reach F} from a note & 42.3\,ms & 6.6\,ms \\
& \texttt{reach F} from a hub & 54.3\,ms & 16.6\,ms \\
& \texttt{walk via F\{1,3\}} from a note & 64.9\,ms & 26.8\,ms \\
\addlinespace
vectors through Arrow and NumPy & \texttt{near} & 54.9\,ms & 19.8\,ms \\
& \texttt{diffuse S} from a note & 58.2\,ms & 23.3\,ms \\
\addlinespace
citation reads batched & \texttt{reach F} from a note & 5.6\,ms & 2.0\,ms \\
& \texttt{reach F} from a hub & 14.3\,ms & 5.4\,ms \\
& \texttt{reach F}, depth 3 & 19.9\,ms & 7.6\,ms \\
& \texttt{walk via F\{1,3\}} from a note & 22.9\,ms & 8.1\,ms \\
& \texttt{walk via F\{1,3\}} from a hub & 19.1\,ms & 7.0\,ms \\
\bottomrule
\end{tabular}
\end{table}

\begin{table}[ht]
\centering\small
\caption{The walk's row, full inverse then component, separate processes; peak resident memory
includes the input matrix.}\label{tab:walk}
\begin{tabular}{@{}lrrrr@{}}
\toprule
Support & Nodes & Full inverse & Component & Peak memory \\
\midrule
fragmented & 300 & 1.7\,ms & 0.5\,ms & 80 then 80\,MB \\
fragmented & 1{,}000 & 21.8\,ms & 3.3\,ms & 160 then 110\,MB \\
fragmented & 2{,}000 & 118\,ms & 12.6\,ms & 440 then 170\,MB \\
fragmented & 5{,}000 & 1{,}228\,ms & 76\,ms & 2.07 then 0.44\,GB \\
fragmented & 10{,}000 & 7{,}252\,ms & 286\,ms & 7.94 then 1.20\,GB \\
connected & 1{,}000 & 18.9\,ms & 24.8\,ms & 160 then 160\,MB \\
connected & 5{,}000 & 1{,}188\,ms & 1{,}351\,ms & 2.07 then 2.27\,GB \\
\bottomrule
\end{tabular}
\end{table}

The fragmented support is a 10-neighbour backbone of clustered nodes, which splits into clusters
as the similarity backbone of clustered notes does; there the component route gains up to 25
times. On a connected support it gains nothing and pays 14\% at 5{,}000 nodes and 31\% at 1{,}000
to find that out. A sentence's support is bounded at 200 nodes per seed, so the gain is for large
supports only.

\begin{table}[ht]
\centering\small
\caption{The spine's partition test by closure and by the index, on synthetic spines of the live
shape.}\label{tab:spine}
\begin{tabular}{@{}rrrrr@{}}
\toprule
Events & By closure & By the index & Every closure decoded & Peak memory \\
\midrule
2{,}801 & 483\,ms & 6.4\,ms & 32.9\,ms & 80 then 60\,MB \\
5{,}601 & 2{,}091\,ms & 14.8\,ms & 123\,ms & 130 then 70\,MB \\
11{,}201 & 8{,}722\,ms & 33.0\,ms & 544\,ms & 460 then 100\,MB \\
22{,}401 & 36{,}484\,ms & 82.4\,ms & 1{,}940\,ms & 1{,}090 then 180\,MB \\
\bottomrule
\end{tabular}
\end{table}

The index is 75 to 443 times faster on the test as written and 15 to 19 times faster even when every
closure is decoded, the cost any consumer that reads all of them must pay. It grows 2.3 times per
doubling of the events against the closure's 4.

\subsection{The verdict}\label{sec:verdict}
\Cref{tab:verdict} applies Proposition~\ref{prop:amdahl} to each candidate. Only the interpreter spends all of
its time in Python, and only there would a compiled rewrite clear a fivefold gain on its own work;
no measured workload spends more than milliseconds in it (a 10{,}000-step script takes 3\,ms, and a
real script's time is in the stores it reads and writes). Batching the meter's charges would take the
quarter of each step the meter costs, in Python.

\begin{table}[ht]
\centering\small
\caption{The candidates, their Python share, and what a compiled rewrite could buy end to end;
where the share is 100\%, the rewrite's own factor, estimated.}\label{tab:verdict}
\begin{tabular}{@{}lrrl@{}}
\toprule
Candidate & Python share & Bound & Outcome \\
\midrule
\texttt{near} family & 10 to 19\% & 1.2$\times$ & no; Arrow read in Python \\
citation-graph family & about 50\% & 2$\times$ & no; batched reads in Python \\
killed-walk solve & under 1\% & none & no; component solve in Python \\
citation search, edge lookup & 4 to 49\% & 1.0 to 1.9$\times$ & no \\
ingress interpreter & 100\% & 20$\times$ or more & yes on its own work; no workload \\
the scorer's bootstrap & 1 to 2\% & none & no \\
spine restriction & 100\% & 10 to 30$\times$ & done in Python, 75 to 443$\times$ \\
\bottomrule
\end{tabular}
\end{table}

A compiled rewrite would cost more than its code. The one pure-Python wheel would become a wheel per
platform and interpreter; the build would move to a Rust-aware backend or gain a second
distribution; a PyO3 release supporting Python 3.15 would be needed, which we did not check; a
toolchain would join the environment's setup; and the import firewall, which reads Python imports,
would not see a crate's dependencies, so the core's ban on networking would need a second
enforcement. Typed stubs would be needed for the type checker. For the interpreter, the gas table,
whose unit is the measured step, would need recalibrating, and the correspondence between the
script laws and their Lean model would need re-establishing. A vectorized rewrite in NumPy or SciPy
costs a function and nothing structural.

\subsection{Re-measurement in a second session}\label{sec:remeasure}
We repeated six figures with scripts of our own on macOS 27.0.1, an Apple M2 Pro with 32\,GiB,
Python 3.15.2, NumPy 2.5.4, SciPy 1.18.1 and pyarrow 26.0.0, importing an export of the branch head
\texttt{6db9f23} (\cref{tab:remeasure}). The HEAD comparison ran on a linked worktree of the
repository, with \texttt{git status} told to take no optional locks, 60 interleaved rounds. The
vector conversion used 304 random vectors of 768 dimensions as an Arrow column. The citation read
used a store with the real schema holding 3{,}000{,}000 edges (60 commits of 50{,}000, a skewed
in-degree) and the 200 most-cited refs at one commit; it times the read primitive alone, both
directions read twice per ref as the builder did, so it exceeds the end-to-end 2.7 times of
\cref{tab:changes}. The walk used 10-nearest-neighbour graphs of points in 16 dimensions, in clusters
of 100 for the fragmented case and one cloud for the connected one; the rows agreed with the full
inverse's to a relative $3 \times 10^{-16}$, with the same unreachable nodes. The reachability check
used DAGs of chains up to 40 long with 0.9 forward cross edges per event and half the events kept,
and the index's result equalled the closure's. The interpreter ran nested loops over a 40-element
list.

\begin{table}[ht]
\centering\small
\caption{Six figures measured again, medians.}\label{tab:remeasure}
\begin{tabular}{@{}lll@{}}
\toprule
Operation & Before & After \\
\midrule
commit of HEAD & 25.8\,ms (two git processes) & 0.22\,ms (files), same sha \\
304 half-precision vectors to lists & 34.2\,ms (row by row) & 2.8\,ms (NumPy), bit for bit \\
304 single-precision vectors to lists & 3.5\,ms & 2.7\,ms, bit for bit \\
200 refs' endpoint pairs & 163.3\,ms (800 statements) & 21.2\,ms (2 statements) \\
walk row, fragmented, 10{,}000 nodes & 6{,}942\,ms, 7.62\,GB & 282\,ms, 0.88\,GB \\
walk row, fragmented, 5{,}000 nodes & 1{,}144\,ms, 1.98\,GB & 71\,ms, 0.27\,GB \\
walk row, connected, 5{,}000 nodes & 1{,}189\,ms, 1.98\,GB & 1{,}326\,ms, 2.18\,GB \\
reachability, 22{,}401 events & 27{,}159\,ms (closure) & 52\,ms (index), 4{,}480\,ms decoded \\
reachability, 2{,}801 events & 337\,ms & 2\,ms, 41\,ms decoded \\
interpreter, 32{,}810 to 524{,}960 steps & & 0.284 to 0.286\,\textmu s a step \\
\bottomrule
\end{tabular}
\end{table}

Three readings follow. The figures agree with the profile in direction and, except where the
fixture differs, in size. The vector conversion's cost was a property of half precision: Arrow's
row-by-row conversion of single-precision vectors is only a quarter slower than the NumPy route,
while that of half-precision vectors is twelve times slower, and the live store is half precision.
And decoding every closure costs more relative to the build on our DAG than on the profile's
spines; the two fixtures differ in shape and in the share of events kept, and we did not separate
the causes. The build itself is two to three orders of magnitude below the closure in both.

\subsection{Packaging checks}
We built the wheel offline from the export, without a repository, so it took the fallback version.
It holds 410 files tagged \texttt{py3-none-any}: the eight packages, the 13 files of the script
library, and the committed defaults (\texttt{defaults.toml}, \texttt{tuning.toml},
\texttt{stack.toml} and one sweep), with no per-machine overlay. Installed into a fresh environment
from the package cache, run from an unrelated directory, it imported \texttt{core},
\texttt{eval.predictions} and \texttt{scheduler}, printed the help of \texttt{python -m
eval.predictions} and of \texttt{mind-palace}; with no root set, the first file system use of a
path under the root raised the plain refusal naming \envroot.

\subsection{Tests and mutations}\label{sec:tests}
The branch's new tests (33, in seven files) and the reference view's 19 pass in 4.2\,s. Each change
keeps the code it replaced as an oracle: a Hypothesis property compares the index's restriction
with the closure's on 300 random spines (events, closed adjacency, generators, order and proper
time); another compares the batched search and wiring with the per-node reads on 200 random edge
sets; the walk's rows are compared with the full inverse at $10^{-10}$; the vector read is
compared float by float, as hexadecimal, with the old read in a half- and a single-precision store;
and the HEAD reader is compared with git in five repository layouts. \Cref{tab:mutations} lists
eight single edits; each fails at least one test. One test, which compares the HEAD reader with git
on the checkout itself, cannot run in an export without a repository and was left out of these runs.

\begin{table}[ht]
\centering\small
\caption{Single mutations, each run against the 51 tests that run in an export.}
\label{tab:mutations}
\begin{tabular}{@{}lrl@{}}
\toprule
Mutation & Fail & Failing tests \\
\midrule
HEAD reader ignores packed refs & 1 & loose, packed and detached refs \\
HEAD reader ignores a worktree's common directory & 1 & the linked worktree \\
vectors narrowed to half precision & 1 & bit-for-bit read, single-precision store \\
batched read returns outgoing edges only & 1 & batched search equals per-node reads \\
the walk solved on node 0's component & 7 & rows equal the full inverse; isolated source \\
index omits the direct successor's bit & 2 & both index-against-closure properties \\
the root variable ignored & 2 & the variable wins; a non-directory refused \\
an unset root usable as a path & 1 & the unresolved root refuses at use \\
\bottomrule
\end{tabular}
\end{table}

\section{Limitations}\label{sec:limits}

\paragraph{Synthetic fixtures.} Every profile figure comes from fixtures of an assumed shape at the
live sizes. The live spine, the citation graph's real degree distribution and a real embedder's
latency were not measured, and a figure within a factor of two of a decision should be re-run on the
live shapes before acting on it.

\paragraph{One machine, live noise.} The profile ran on one laptop with the instance's daemon
running, and the re-measurement on one machine. Medians of a few runs reduce but do not remove
setup effects of the kind \citet{mytkowicz2009} describe.

\paragraph{An uncommitted harness.} The profile's harness was not committed. Its parameters are
stated and six of its figures were reproduced independently, but the others cannot be regenerated
exactly from the repository.

\paragraph{Estimates in the verdict.} The interpreter's gain from a compiled rewrite and the spine's
constant-factor gain are estimates, not measurements, and are marked as such.

\paragraph{Where the changes do not help.} The component solve gains nothing on a connected
support and costs 14 to 31\%; a giant component stays cubic. The index's memory is $n^2/64$ words,
and any consumer that reads every closure still pays for a quadratic output. The HEAD reader declines
an unborn branch or a reftable repository, which then costs the git processes again.

\paragraph{Packaging gaps.} A top-level package named \texttt{eval} claims that name in an
environment's site packages. The version falls back to \texttt{0.0.0+unknown} without a repository,
so a build from a source tarball carries no version.

\section{Future work}
Committing the harness as a benchmark suite would let every figure here be regenerated, and running
it in CI would turn the verdict into a regression gate. Re-running the profile on the live shapes,
with a real embedder, would test the fixtures' assumptions. Charging the interpreter's meter in
batches would take a quarter of each step. A giant connected support remains cubic; an iterative or
low-rank approximation of the needed rows would change that, at the price of the exactness the walk's
laws are stated for. A source-distribution check in CI would catch a build that loses its version.

\section{Conclusion}
Packaging needed one rule for where an instance's state lives and a refusal that waits until the
state is used. Measuring before rewriting, with controls for what the profiler cannot see, showed
that the candidates' time was mostly in compiled libraries already, or in work that did not need
doing: a git call that two files answer, a conversion done row by row, statements issued per node, a
solve over components that do not matter, and a closure stored when a mask would do. Five changes in
Python removed most of it, each pinned by a test against the code it replaced. The one operation that
is entirely Python, the interpreter, has no workload that would feel a compiled rewrite, so the
codebase stays a single pure-Python wheel.

\bibliographystyle{plainnat}
\bibliography{refs}
\end{document}
